%% kapa: exam=1 round=19 year=2008 period=1 form=A subject=civil_law subject_ko=민법 seq=19 number=19
%% source: Q-Net 감정평가사 기출문제 (cst00309) artlSeq=1001962 file=extracted/1008/1001962/2008년 제19회 감정평가사 1차시험 기출문제/1교시.hwp
%% conversion: kordoc 4.24.51 -> tools/gen1.py (auto); review=proofread-sonnet; flags=none
\begin{question}{19}
대리권의 범위와 제한에 관한 설명으로 틀린 것은?(다툼이 있는 경우에는 판례에 의함)
\begin{choices}
  \item 부동산 경매절차에서 동일물건에 관하여 이해관계가 다른 2인 이상의 대리인이 된 경우에 그 대리인이 한 매수신청(입찰)은 유효하다.
  \item 부동산을 매도할 권한이 있는 대리인은 특별한 사정이 없는 한 그 잔대금도 수령할 권한이 있다.
  \item 매매계약의 체결과 이행에 관하여 포괄적으로 대리권을 수여받은 대리인은 특별한 다른 사정이 없는 한 상대방에 대하여 약정된 매매대금지급기일을 연기하여 줄 권한도 가진다.
  \item 특별한 사정이 없는 한 부동산을 매수할 권한을 수여받은 대리인에게 당연히 그 부동산을 처분할 권한이 주어지는 것은 아니다.
  \item 법정대리인인 친권자가 자신의 부동산을 미성년자인 자녀에게 증여하고자 자녀를 대리하여 계약을 체결하는 것은 자기계약이지만 유효하다.
\end{choices}
\end{question}
\answer{1}
